\chapter{平方根}\label{ch:g9:sqrt}

边长为 $1$ 的正方形的对角线有多长？毕达哥拉斯定理
给出 $\sqrt 2$，它是平方为 $2$ 的数，
而它并不是分数。本章定义平方根，建立
平方根的运算法则，并讲解如何化简
形如 $\sqrt{75}$ 的式子。

\section{定义与基本性质}

\begin{definition}[平方根]\label{def:g9:sqrt:def}
设 $a \geq 0$。$a$ 的\emph{平方根}\index{平方根}记作
$\sqrt a$，它是平方为 $a$ 的唯一\emph{非负}数：
\[
\sqrt a \geq 0
\qquad\text{且}\qquad
\left(\sqrt a\right)^2 = a .
\]
负数没有平方根，因为任一数的平方都非负。
\end{definition}

\begin{example}
$\sqrt{49} = 7$, $\sqrt 0 = 0$, $\sqrt 1 = 1$,
$\sqrt{2.25} = 1.5$。需要熟记的第一批平方根是
\emph{完全平方数}的平方根：$1, 4, 9, 16, 25, 36, 49, 64, 81, 100, 121,
144$。
\end{example}

\begin{omfigure}
\begin{tikzpicture}[scale=1.35]
  \draw[thick, fill=omDef!8] (0,0) rectangle (2,2);
  \draw[omThm, thick] (0,0) -- (2,2);
  \node[below, font=\small] at (1,0) {$1$};
  \node[left, font=\small] at (0,1) {$1$};
  \node[omThm, above left=-2pt, font=\small, rotate=45] at (1.08,0.92)
    {$\sqrt2$};
  \draw[thin] (1.75,0) -- (1.75,0.25) -- (2,0.25);
\end{tikzpicture}

{\small 单位正方形的对角线长为 $\sqrt{1^2 + 1^2} =
\sqrt2 \approx 1.414$，由毕达哥拉斯定理
（\cref{ch:g9:trig}中已回顾此定理）。这个数是无理数：它的小数
永不循环。}
\end{omfigure}

\begin{proposition}[平方数的平方根]\label{prop:g9:sqrt:square}
对任一实数 $x$（不论正负）：
\[
\sqrt{x^2} = \abs{x} =
\begin{cases}
x & \text{若 } x \geq 0,\\
-x & \text{若 } x < 0 .
\end{cases}
\]
\end{proposition}

\begin{proof}
$\abs{x}$ 非负，且其平方为 $x^2$（一个数与其相反数的
平方相等）。由于它是平方为
$x^2$ 的唯一非负数，故它即为 $\sqrt{x^2}$。
\end{proof}

\begin{example}
$\sqrt{(-5)^2} = \sqrt{25} = 5 = \abs{-5}$：平方根“忘记”了符号，
它并不会恢复符号。特别地，$\sqrt{x^2} = x$ 对
负数 $x$ 是\emph{错误}的。
\end{example}

\section{平方根的乘积与商}

\begin{theorem}[乘法与除法法则]\label{thm:g9:sqrt:rules}
对所有 $a \geq 0$ 与 $b \geq 0$：
\[
\sqrt{a b} = \sqrt a \times \sqrt b,
\qquad\text{且对 } b > 0:\quad
\sqrt{\frac ab} = \frac{\sqrt a}{\sqrt b} .
\]
\end{theorem}

\begin{proof}
$\sqrt a \times \sqrt b$ 非负（非负数之积），
且其平方为
\[
\left(\sqrt a \times \sqrt b\right)^2
= \left(\sqrt a\right)^2 \times \left(\sqrt b\right)^2
= a b .
\]
由平方为 $ab$ 的非负数唯一，它等于
$\sqrt{ab}$。商的法则同理可证。
\end{proof}

\begin{remark}[和的情形没有此法则！]
一般地，$\sqrt{a + b}$ \emph{不}等于 $\sqrt a + \sqrt b$：
\[
\sqrt{9 + 16} = \sqrt{25} = 5,
\qquad\text{而}\qquad
\sqrt 9 + \sqrt{16} = 3 + 4 = 7 .
\]
\end{remark}

\begin{method}[化简 $\sqrt n$]\label{met:g9:sqrt:simplify}
要化简整数的平方根：
\begin{enumerate}
  \item 求能整除 $n$ 的最大完全平方数（必要时
        先分解 $n$）；
  \item 利用乘积法则拆分，并开出平方：
        $\sqrt{k^2 m} = k\sqrt m$；
  \item 检验根号内剩下的
        部分中不再含完全平方数因子。
\end{enumerate}
\end{method}

\begin{example}\label{ex:g9:sqrt:simplify}
化简 $\sqrt{75}$：因为 $75 = 25 \times 3$，
\[
\sqrt{75} = \sqrt{25 \times 3} = \sqrt{25} \times \sqrt 3 = 5\sqrt3 .
\]
化简 $\sqrt{72}$：因为 $72 = 36 \times 2$，
$\sqrt{72} = 6\sqrt2$。再看一个商：
$\sqrt{\dfrac{49}{16}} = \dfrac{\sqrt{49}}{\sqrt{16}} = \dfrac74$。
\end{example}

\begin{example}[根式的加法]\label{ex:g9:sqrt:add}
只有当各根式\emph{同类}时，根式的和才能化简。
计算 $\sqrt{12} + \sqrt{27} - \sqrt{48}$，先化简每一项：
\begin{align*}
\sqrt{12} &= \sqrt{4 \times 3} = 2\sqrt3, \\
\sqrt{27} &= \sqrt{9 \times 3} = 3\sqrt3, \\
\sqrt{48} &= \sqrt{16 \times 3} = 4\sqrt3,
\end{align*}
于是和为 $2\sqrt3 + 3\sqrt3 - 4\sqrt3 = (2 + 3 - 4)\sqrt3 = \sqrt3$。
\end{example}

\begin{example}[根式的展开]\label{ex:g9:sqrt:expand}
代数中的恒等式同样适用于根式。展开
$\left(\sqrt5 + 2\right)^2$，用 $(a+b)^2 = a^2 + 2ab + b^2$：
\[
\left(\sqrt5 + 2\right)^2 = 5 + 2 \times 2\sqrt5 + 4 = 9 + 4\sqrt5 .
\]
再用第三个恒等式 $(a+b)(a-b) = a^2 - b^2$：
\[
\left(\sqrt7 + \sqrt3\right)\left(\sqrt7 - \sqrt3\right) = 7 - 3 = 4 :
\]
两个无理数的积可以是整数。
\end{example}

\section{方程 \texorpdfstring{$x^2 = a$}{x2 = a}}

\begin{theorem}[解方程 $x^2 = a$]\label{thm:g9:sqrt:equation}
\begin{enumerate}
  \item 若 $a > 0$，则方程 $x^2 = a$ 恰有两个解：
        $\sqrt a$ 和 $-\sqrt a$；
  \item 若 $a = 0$，则唯一解为 $0$；
  \item 若 $a < 0$，则方程无解。
\end{enumerate}
\end{theorem}

\begin{proof}
将 $x^2 = a$ 改写为 $x^2 - \left(\sqrt a\right)^2 = 0$，
在 $a \geq 0$ 时，利用平方差公式分解因式：
\[
\left(x - \sqrt a\right)\left(x + \sqrt a\right) = 0 ,
\]
于是 $x = \sqrt a$ 或 $x = -\sqrt a$（$a = 0$ 时两者重合）。当
$a < 0$ 时方程无解，因为 $x^2 \geq 0 > a$ 对所有 $x$ 成立。
\end{proof}

\begin{omfigure}
\begin{tikzpicture}
\begin{axis}[omaxis,
  width=7.6cm, height=5.4cm,
  xmin=-3.1, xmax=3.1, ymin=-1.6, ymax=5.4,
  xtick={-1.732,1.732}, xticklabels={$-\sqrt3$,$\sqrt3$},
  ytick={3}, yticklabels={$3$}]
  \addplot[omDef, thick, domain=-2.25:2.25, samples=90] {x^2};
  \addplot[omThm, thick, domain=-2.9:2.9] {3};
  \fill (axis cs:-1.732,3) circle (1.5pt);
  \fill (axis cs:1.732,3) circle (1.5pt);
  \draw[dashed] (axis cs:-1.732,3) -- (axis cs:-1.732,0);
  \draw[dashed] (axis cs:1.732,3) -- (axis cs:1.732,0);
\end{axis}
\end{tikzpicture}

{\small 求解 $x^2 = 3$，即抛物线 $y = x^2$ 与水平直线
$y = 3$ 相交于横坐标为 $\pm\sqrt3$ 的两点。}
\end{omfigure}

\begin{example}
$x^2 = 16$：解为 $4$ 与 $-4$。\quad
$x^2 = 5$：解为 $\sqrt5$ 与 $-\sqrt5$（精确值；$\sqrt5
\approx 2.236$）。\quad
$x^2 = -9$：无解。\quad
$3x^2 = 21$：先两边同除以 3，得 $x^2 = 7$，解为 $\pm\sqrt7$。
\end{example}

\section{练习}

\begin{exercise}[$\star$]\label{exo:g9:sqrt:1}
不用计算器计算：
\[
\sqrt{64}, \qquad \sqrt{100}, \qquad \sqrt{0.09}, \qquad
\sqrt{\tfrac{1}{25}}, \qquad \left(\sqrt{13}\right)^2, \qquad
\sqrt{(-4)^2} .
\]
\end{exercise}

\begin{exercise}[$\star$]\label{exo:g9:sqrt:2}
化简：
\[
\sqrt{50}, \qquad \sqrt{45}, \qquad \sqrt{98}, \qquad \sqrt{300} .
\]
\end{exercise}

\begin{exercise}[$\star$]\label{exo:g9:sqrt:3}
计算并化简：
\[
\sqrt2 \times \sqrt{18}, \qquad
\frac{\sqrt{75}}{\sqrt3}, \qquad
\sqrt5 \times \sqrt{20}, \qquad
\frac{\sqrt{8}}{\sqrt{50}} .
\]
\end{exercise}

\begin{exercise}[$\star$]\label{exo:g9:sqrt:4}
解下列方程：$x^2 = 36$；\quad $x^2 = 11$；\quad $x^2 + 4 = 0$；\quad
$2x^2 = 50$。
\end{exercise}

\begin{exercise}[$\star$]\label{exo:g9:sqrt:5}
合并为一项：$\sqrt{20} + \sqrt{45}$，然后
化简 $3\sqrt8 - \sqrt{18} + \sqrt2$。
\end{exercise}

\begin{exercise}[$\star\star$]\label{exo:g9:sqrt:6}
展开并化简：
\[
\left(\sqrt3 + 1\right)^2, \qquad
\left(2\sqrt5 - 3\right)^2, \qquad
\left(\sqrt6 + \sqrt2\right)\left(\sqrt6 - \sqrt2\right).
\]
\end{exercise}

\begin{exercise}[$\star\star$]\label{exo:g9:sqrt:7}
一块正方形场地的面积为 $6400$ m$^2$。它的边长是多少？
其对角线长是多少（先给出精确值，再四舍五入到整数米）？
\end{exercise}

\begin{exercise}[$\star\star$]\label{exo:g9:sqrt:8}
证明 $\dfrac{1}{\sqrt2} = \dfrac{\sqrt2}{2}$（分子与分母同乘
$\sqrt2$）。用同样的方法，将 $\dfrac{6}{\sqrt3}$
和 $\dfrac{10}{\sqrt5}$ 化为分母不含根号的形式。
\end{exercise}

\begin{exercise}[$\star\star$]\label{exo:g9:sqrt:9}
判断下列命题的真假，并用证明或反例说明理由。
\begin{enumerate}
  \item 对所有 $a, b \geq 0$：$\sqrt{ab} = \sqrt a \sqrt b$。
  \item 对所有 $a, b \geq 0$：$\sqrt{a+b} = \sqrt a + \sqrt b$。
  \item 对所有 $x$：$\sqrt{x^2} = x$。
  \item $\left(3\sqrt2\right)^2 = 18$.
\end{enumerate}
\end{exercise}

\begin{exercise}[$\star\star\star$]\label{exo:g9:sqrt:10}
设 $x = \sqrt{7 + 4\sqrt3}$。计算
$\left(2 + \sqrt3\right)^2$，并由此求出 $x$ 的更简单的表达式。
对 $\sqrt{9 - 4\sqrt5}$ 同样处理（目标是化为
$\left(\sqrt5 - 2\right)^2$ 的形式，并注意符号）。
\end{exercise}

\section{问题：不是分数的数}

\begin{problem}[{周末问题 --- 传奇的证明，
$\sqrt2$ 是无理数，它的众多“近亲”，以及海伦精妙的
近似方法}]
\label{pb:g9:sqrt:1}
单位正方形的对角线是一个完全实在的长度 --- 你
在\cref{ex:g8:pythagoras:diagonal}中亲手画过它 --- 然而，正如
毕达哥拉斯学派在约两千五百年前惊恐地发现，
\emph{任何分数都无法度量它}。相传
这一发现曾招致溺死的惩罚。本题带你完成那个不朽的证明
（四个问题，终身受用），
再让众多“受害者”现身，并以
工程师两千年间用来驯服这个不可驯服之
数的方法作结。

\medskip
\noindent\textbf{第一部分 --- 证明。}
\begin{enumerate}
  \item 先把它逼近：验证
        $1.4 < \sqrt2 < 1.5$，再验证 $1.41 < \sqrt2 < 1.42$，
        将各界平方即可核实。再多取一位：
        $\sqrt2$ 介于哪两个三位小数之间？
  \item 现在假设 --- 为了导出矛盾 --- 存在
        使 $\sqrt2 = \frac ab$ 成立的\emph{既约分数}
        （\cref{met:g9:arith:simplify}）。两边平方，
        并证明 $a^2 = 2 b^2$。
        $a^2$ 的奇偶性如何？
  \item 由\cref{pb:g8:pythagoras:1}的奇偶性事实可知：奇数
        的平方为奇数。由此推出 $a$ 为偶数，
        设 $a = 2k$ 并代入 --- 证明 $b$ 必须
        也为偶数。
  \item 矛盾在哪里？由此得出结论，并完整陈述
        定理：\emph{$\sqrt2$ 是无理数 --- 它不是
        任何两个整数的商}。
  \item 证明曾一度不动声色地依赖“既约”这几个字。
        指出缺了它则哪一步会失效，
        并说明为何一开始可以假设分数是既约的，
        这一假设本身是合理的。
\end{enumerate}

\medskip
\noindent\textbf{第二部分 --- “受害者”成倍增加。}
\begin{enumerate}[resume]
  \item 第二种证明，借助质因数分解
        （\cref{thm:g9:arith:factorization}）：在 $a^2 = 3b^2$ 中，
        比较两边的\emph{$3$ 的指数}的奇偶性
        （见\cref{pb:g9:arith:1}：平方数的
        质因数指数均为偶数）。由此得出 $\sqrt3$ 为无理数。
  \item 对 $a^2 = n b^2$ 作同样的指数论证，其中
        $n$ 为一般的整数：对哪些 $n$ 能导出矛盾，
        哪些则会失效？陈述完整结论：
        $\sqrt n$ 为无理数，当且
        仅当\,\dots
  \item 证明“盾引理”：非零有理数乘无理数
        仍是无理数。（设 $r x = q$，其中 $r, q$
        为有理数，$r \neq 0$，并解出 $x$。）由此推出
        $2\sqrt2$ 与 $\frac{\sqrt2}{2}$ 均为无理数。
  \item 证明 $1 + \sqrt2$ 为无理数。再看更漂亮的一题：
        假设 $s = \sqrt2 + \sqrt3$ 为有理数，
        计算 $s^2$，分离出 $\sqrt6$，并找出
        矛盾。
  \item 稍作冷静：给出两个无理数，它们的
        \emph{和}为有理数，再给出两个无理数，它们的\emph{积}为
        有理数。（无理性并不在四则运算下保持
        --- 每种情形都需要单独证明。）
\end{enumerate}

\medskip
\noindent\textbf{第 III 部分 --- 海伦的算法。}
早在计算器出现的两千年前，亚历山大的海伦
用如下方法逼近 $\sqrt2$：\emph{猜测 $x$；
用 $x$ 与 $\frac2x$ 的平均数代替猜测值；重复。}
\begin{enumerate}[resume]
  \item 从猜测值 $x = 1$ 出发，
        依次计算接下来的两个猜测值，并用精确分数表示。
  \item 计算第三个猜测值，仍用精确分数表示，
        并求出其小数值。与第 1 题相比，
        $\sqrt2$ 已有多少位小数是正确的？
  \item 该算法背后的思想是：面积为 $2$ 的长方形，
        若一边为 $x$，则另一边为 $\frac2x$。证明
        $\sqrt2$ 总\emph{介于} $x$ 与 $\frac2x$ 之间
        （考虑两种情形 $x^2 > 2$ 与 $x^2 < 2$），因此
        对两边取平均，即可使长方形
        向正方形逼近。
  \item 猜测值 $\frac32$、$\frac{17}{12}$、
        $\frac{577}{408}$ 中藏着一个妙处：计算 $3^2 - 2 \times
        2^2$，再计算 $17^2 - 2 \times 12^2$，再计算
        $577^2 - 2 \times 408^2$。第 I 部分已证明
        $a^2 - 2b^2 = 0$ 不可能成立 --- 海伦的分数
        与这种不可能情形有多接近？
  \item 结语，用两句话概括：单位正方形的对角线在纸上存在，
        但没有任何分数能够度量它，其小数写法
        永远不会循环
        （见\cref{pb:g9:fractions:1}）。那么数轴
        必须包含哪一类“新数”，它们的完整理论
        在本系列的何处讲述？
\end{enumerate}
\end{problem}
